\section*{The object being priced}

A stablecoin basket is a portfolio of tokens each of which claims to be worth one dollar. Basket margin at centralised venues, index perpetuals on decentralised ones and proof-of-reserves programmes at regulated issuers all price the same thing: the joint behaviour of the basket when the claim fails for one of its members. What they need is the joint distribution of peg deviations, and in practice what they use is its second moment, the correlation matrix.

That matrix is not observed. On a representative day only three or four coins in a panel of dozens deviate from par at all, and the rest supply a flat line of near-zero returns. A correlation computed on such data is dominated by that flatness, and the matrices that clearinghouses, index providers and regulators publish are consequently smoothed across the gaps where the deviations are missing. The smoothing is usually defended as interpolation: with only a handful of non-trivial observations, some regularisation is unavoidable.

\section*{Why the gaps are structural}

We argue that this defence is misplaced, because the gaps are not sampling artefacts. They are regime boundaries. A stablecoin's peg survives or fails according to structural conditions that change discretely: whether its collateral backs it at par \citep{LyonsViswanath2023}, whether its price feed is still reporting \citep{LiuLiSprenger2023}, whether redemption is open or gated, whether the custodian holding its reserves is solvent \citep{MaZengZhang2023}, whether a bridge has been compromised, whether a regulator has intervened, whether the basis on the funding leg has compressed, and whether a PSM buffer still absorbs the spread. Each of these is a channel through which a basket can be stressed, and each has a threshold.

Within a regime - a region of the parameter space in which no channel has crossed its threshold - the joint co-movement structure of the basket does not change. It is not that it changes slowly; it does not change. Across a threshold it steps. Figure~\ref{fig:wall} shows the geometry for the two channels that admit a continuous representation, collateral backing and price-feed deviation, and locates the remaining six as discrete triggers.

The consequence for correlation modelling is direct. A smooth surface fitted to within-regime data has no way to represent the step, because the step is not in the data it is fitted to: within a regime there is nothing to fit. No amount of additional within-regime data improves the estimate at the boundary, and a smoother that is continuous in the parameters will place at the boundary a value that the true correlation never takes.

\section*{What we posit, and what follows}

We make three assumptions, and we state them as assumptions rather than deriving them, so that a reader can separate what the paper guarantees from what it claims.

The first is that within a regime the co-movement covariance of the basket is constant up to a positive scale. The scale is allowed to move, because the general level of peg-deviation volatility plainly does move; what is fixed is the proportionality structure. We define the covariance as a centred second moment of the peg-deviation vector under the probability law of that regime, which is the object that a correlation estimated from a window of prices actually recovers.

The second is that across a regime boundary the covariance jumps by a signed sum of rank-one terms. Each term is carried by a direction supported on the coins that the crossing channel affects, so the jump has a block structure: the coins a channel touches move together in a way the rest of the basket does not share. The third assumption is that within an episode the support of the jump is a single group of coins, the wall-affected block, fixed by the channel and not by the data.

Everything downstream of these three is proved. The covariance is stated where the crossing law is admissible, because an additive rank-one jump supported on two or more coins necessarily adds a positive amount to its own diagonal, and a correlation matrix cannot absorb a diagonal increment - its diagonal is one on both sides of the boundary by construction. Passing to correlations through the normalisation map removes the obstruction exactly, because that map assigns the unit diagonal rather than computing it. What survives the passage is not the rank-one structure: the normalising map is a nonlinear congruence, and it destroys that. What survives is a signed footprint in three tiers, positive on pairs whose two coins both lie in the block, of opposite sign on pairs with one coin in and one out, and zero on pairs outside the block. And what follows for any smooth estimator is a gap: a locally continuous surface cannot track a function that is constant on regimes and steps between them, so its distance from the true correlation is of the order of the jump and does not shrink as the boundary is approached.

The same reasoning fixes the sign of the block entries. Differentiating the induced correlation jump with respect to the crossing amplitude shows that a block pair rises at the boundary provided its pre-event covariance lies below the harmonic mean of the two variances, a condition that holds automatically when the two coins' variances are equal and that can fail only for a pair that is both strongly correlated and strongly heteroscedastic before the event.

\section*{Testing the block signature}

The three-tier footprint suggests a test that does not require the crossing amplitude, the regime parameters, or any distributional assumption. Classify every pair of coins in an episode by how many of its members lie in the wall-affected block, and ask whether the change in correlation from the pre-event window to the acute window is largest on the block layer. Because a crisis raises co-movement across the whole panel and not only inside the block, the comparison that carries the inference is the block-versus-rest contrast, which differences out a component common to all pairs.

Two details of the inference matter. Pairs share coins and so are not independent, which makes a permutation over pairs anti-conservative; the test we report redraws instead which \emph{coins} form a block of the same size, which asks how unusual the true block is among all the coin sets that could have been so labelled. And because the block is fixed by the channel rather than estimated, the design is a test rather than a fit.

We apply the test to the nine documented de-peg episodes of 2020 to 2025 (Table~\ref{tab:events}), drawn entirely from public price archives. Three of the nine carry a block of more than one coin; the other six involve a single coin and cannot generate a block pair at all.

On the custodial failure of March 2023, in which the reserve custodian of the largest dollar stablecoin failed over a weekend, the block is the three coins exposed through that custodian and through the collateral channel that transmitted the shock: the two largest decentralised stablecoins held the affected token as collateral or as backing. All three layers are ordered as the footprint requires, the block-versus-rest contrast is 0.69, and the coin-level permutation places the true block at the 98.9th percentile of the 816 coin sets of size three that could have been labelled (p = 0.011) (Fig.~\ref{fig:gradient}, Table~\ref{tab:gradient}). The ordering holds across all nine specifications of a window grid that varies the pre-event window between 30 and 90 days and the acute window between three and ten days either side.

The result most exposed to a trivial explanation is the pair formed by the two decentralised coins, which contains neither the token whose custodian failed nor any coin that held reserves at that custodian directly. A victim-only account predicts no jump for this pair. Its members are in the block because the collateral channel touches them, so the test does not establish an effect free of any common exposure; what it establishes is that the effect does not require the coin that actually broke. Its jump, 1.18, is the largest of any pair in the panel.

\section*{Scope of the evidence}

We report the limits as part of the result. Of the two remaining events with a multi-coin block, one is consistent in sign but imprecise, its block-versus-rest contrast positive in eight of nine window specifications, and the other shows no block signal under the canonical window and turns positive only under a 30-day pre-event window. The test therefore finds a clear effect on one episode, a consistent but imprecise sign on a second, and nothing on a third. We do not pool them into a claim about de-peg episodes in general.

Two further conditions bound what the exercise establishes. The block-versus-rest contrast removes a crisis-wide component common to all pairs only if that component is additive and does not differ systematically by layer; we adopt that as an identifying condition rather than as a consequence of the theory, and we report the level of the outside layer so that a reader can see how much of the raw co-movement the subtraction absorbs. And the wall assignment in the historical record draws on post-mortem accounts. In live use the channel is identified by the failure event itself, which is observable as it happens; the instrument described here is therefore a monitor of a channel's incidence, not a forecast of its onset.

\section*{What a supervisor can do with it}

The practical content of the result is narrow and, we think, useful. A supervisor holding only public price data can compute, for any episode, how much more the co-movement of the coins a channel touches has risen than that of the rest of the basket. That statistic requires neither proprietary data nor issuer disclosure, neither of which is available in this market: stablecoin issuers do not publish liquid equity prices and do not disclose bilateral exposures, which is why equity-based systemic-risk measures have found little traction here.

Two uses follow. The first is diagnostic: a supervisor observing a custodial or collateral failure can tell within days whether the basket is transmitting through the channel or merely co-moving with the market, because the two produce different footprints across pair layers. The second is in calibration: a margin or haircut rule that prices the joint tail on a smooth correlation surface will understate it at exactly the moments the surface is asked to carry weight. The size of that understatement is not identified by anything in this paper - the scale of the crossing is left free - but its sign and its support are, and for a rule that needs to know which coins to haircut together rather than by how much, that is the binding input.

The wider point is that stablecoin basket risk is not a smooth function of market conditions but a staircase. Between boundaries nothing happens; at a boundary the structure changes and does not change back until another boundary is crossed. A risk framework that prices the slope is pricing something that, on this evidence, is not there.
